% IV. Методика и резултати. Обединява предишните раздели V и VI.
% ВНИМАНИЕ: тук не се пише нито едно число, преди то да е измерено.
\section{Методика и резултати}
\label{sec:method}
\label{sec:results}

\subsection{Обхват и условия на измерването}
\label{sec:method-scope}

Този параграф определя какво изобщо може да се твърди. \textbf{Разгърнато е при един доставчик}:
описанието на инфраструктурата е приложено в Azure и услугата е измерена от външен клиент
(Раздел~\ref{sec:azure-run}). \textbf{При AWS не е разгръщано}: за тази сметка няма кредит.
Следователно число за закъснение, мащабиране или цена при AWS не съществува и не се съобщава, а
сравнението между двамата доставчици липсва. Такова число не е и оценявано по аналогия:
правдоподобната стойност е по-лоша от липсващата, защото изглежда като измерване. \textbf{Измерено е и локално}: три експеримента върху стек от контейнери на една машина.
Те отговарят не на въпроса за разликата между доставчиците, а на друг: работи ли построената
система, как се държи капацитетът ѝ при промяна на броя реплики, и реагира ли контролерът в разумно
време.

Всички резултати са получени при Windows 11 Pro N (10.0.26200), 12 логически процесора,
g++ 15.2.0 (MinGW-w64) в режим Release, CMake 4.3.2 с Ninja 1.13.2, Docker 29.4.3, образ върху
Alpine Linux 3.20, изпълняван от \code{scratch}. Генераторът се изпълнява на \emph{същата} машина,
върху която работят контейнерите, и дели с тях същите 12 процесора. Това е основният източник на
разсейване в локалните резултати.

Единицата работа е фиксирана за всички измервания: $n = 384$, $m = 500$. Всяка реплика има една
нишка за обработка, така че капацитетът на набора е точно равен на броя реплики. Всяко измерване
има 5 s загряване, чиито резултати се отхвърлят, и 20 s измерване, и се повтаря 3 пъти. Нищо не се
осреднява преди да достигне до файл: съобщава се медианата, но и трите стойности са в
\code{results/}.

\subsection{Капацитет спрямо брой реплики}

Броят реплики се задава на 1, 2, 4 и 6, при затворен контур с 2 едновременни заявки на реплика.
Затвореният контур е избран нарочно: при него отчетената пропускателна способност е мярка за
капацитет, а не мярка за търпението на клиента. Разсейването между повторенията е показано, защото
е част от резултата \tabref{tab:throughput}. Неуспешни заявки няма в дванадесетте повторения.

\begin{table}[H]
\caption{Медиана от три повторения, с най-малката и най-голямата измерена стойност}
\label{tab:throughput}
\centering
\small
\begin{tabular}{@{}rrrrrrr@{}}
\toprule
\textbf{Реплики} & \textbf{Пропуск.} & \textbf{Най-малка} & \textbf{Най-голяма} &
\textbf{p50} & \textbf{p90} & \textbf{p99} \\
 & [зая./s] & [зая./s] & [зая./s] & [ms] & [ms] & [ms] \\
\midrule
1 & 18{,}90 & 16{,}75 & 19{,}80 & 99{,}25  & 115{,}88 & 203{,}35 \\
2 & 33{,}45 & 21{,}45 & 34{,}10 & 129{,}18 & 183{,}66 & 285{,}80 \\
4 & 57{,}95 & 35{,}15 & 59{,}90 & 118{,}35 & 232{,}10 & 462{,}00 \\
6 & 65{,}10 & 38{,}90 & 77{,}60 & 120{,}22 & 390{,}38 & 863{,}21 \\
\bottomrule
\end{tabular}
\end{table}

\textbf{Коментар.} Мащабирането е подлинейно и се изчерпва. От 1 до 4 реплики ускорението е
3{,}07 при четирикратен ресурс, тоест ефективност около 77\%. От 4 до 6 реплики пропускателната
способност нараства с 12\%, а закъснението на 99-и персентил се удвоява, от 462 ms на 863 ms.
Обяснението не е в услугата: при 6 реплики върху 12 логически процесора едновременно работят 6
изчислителни нишки, 64 нишки на посредника и 12 нишки на генератора, тоест системата е презаписана
и се измерва планировчикът на операционната система. Разсейването расте заедно с броя реплики,
затова локалните числа са характеристика на конкретната машина, а не на услугата. Медианното
закъснение остава между 99 и 130 ms, докато високите персентили растат многократно: средната
стойност при 6 реплики би скрила 863 ms опашка зад 120 ms медиана.

\subsection{Време за реакция при увеличаване}

Наборът се увеличава от 1 на 2 реплики и се измерва времето от подаването на командата до момента,
в който посредникът започва да насочва заявки към новата реплика. При 5 повторения измерените
стойности са 4{,}662, 5{,}234, 1{,}357, 2{,}001 и 1{,}840 s, тоест медиана 2{,}001 s. Времето
съдържа стартиране на контейнера, регистрация на адреса във вградения DNS и откриването му от
посредника, чийто период на опресняване е 1 секунда. Разликата от близо четири пъти при иначе
еднакви условия показва, че политика с период на управление от порядъка на секунда работи с шум от
същия порядък. Затова периодът в следващия експеримент е 3 секунди.

\subsection{Поведение на контролера под натоварване}

Стекът е вдигнат с една реплика, контролерът е стартиран с цел 0{,}70, граници от 1 до 6 реплики,
период 3 s и изчакване 6 s при увеличаване и 15 s при намаляване. Натоварването е в отворен контур,
40 заявки в секунда, 32 едновременни, 45 s измерване. Отчетът на генератора:
1868 успешни заявки, 0 неуспешни, 41{,}51 заявки/s, p50 = 51{,}84 ms, p90 = 1508{,}20 ms,
p99 = 1639{,}47 ms. Високите персентили са опашката, натрупана в първите секунди, когато една
реплика приема натоварване, което не може да поеме.

\begin{table}[H]
\caption{Извадка от решенията на контролера. Пълната редица е в \code{results/scaling-timeline.csv}}
\label{tab:timeline}
\centering
\small
\begin{tabular}{@{}rrrrL{5.0cm}@{}}
\toprule
\textbf{$t$ [s]} & \textbf{Реплики} & \textbf{Натоварване} & \textbf{Желани} & \textbf{Решение} \\
\midrule
6{,}9  & 1 & 0{,}000 & 1 & задържане, на целта \\
10{,}3 & 1 & 1{,}101 & 2 & \textbf{увеличаване}, над целта \\
15{,}6 & 2 & 0{,}848 & 2 & задържане, изчакване преди увеличаване \\
20{,}2 & 2 & 0{,}822 & 3 & \textbf{увеличаване}, над целта \\
25{,}6 & 3 & 0{,}996 & 3 & задържане, изчакване преди увеличаване \\
33{,}3 & 3 & 0{,}667 & 3 & задържане, на целта \\
57{,}9 & 3 & 0{,}646 & 3 & задържане, на целта (десети пореден) \\
60{,}9 & 3 & 0{,}033 & 1 & \textbf{намаляване}, под целта \\
65{,}9 & 1 & \multicolumn{2}{l}{нулиран брояч} & пропуснат интервал \\
\bottomrule
\end{tabular}
\end{table}

\textbf{Коментар.} Контурът се затваря \tabref{tab:timeline}. Натоварването расте, контролерът
увеличава набора на две стъпки, след което то се установява под целта 0{,}70 и решенията спират за
десет интервала. След премахването на натоварването наборът се връща на една реплика. Установената стойност позволява независима проверка на цялата верига: при 3
реплики, натоварване 0{,}655 и 40{,}0 заявки в секунда времето за обслужване на заявка следва да е
$0{,}655 \cdot 3 / 40{,}0 = 49{,}1$ ms, което съвпада с времето, което изчислителното ядро отчита
само за себе си. Тоест метриката измерва извършената работа, а не чакането до нея.

Редовете при $t = 15{,}6$ и $t = 25{,}6$ показват периода на изчакване: без него контролерът щеше
да добавя реплики, докато първите се стартират.

\subsection{Разгръщане в Azure Container Apps}
\label{sec:azure-run}

Средата на разгръщането е следната. Регион \code{italynorth} (Италия Север, Милано). Адрес
\code{stratus-dev.agreeablemeadow-0a00d2a7.italynorth.azurecontainerapps.io}. Образ
\code{stratusacr4380.azurecr.io/stratus:v1}. Контейнер с 1 vCPU и 2{,}00 GiB памет, при
\code{STRATUS\_THREADS=1}, тоест една изчислителна нишка на реплика. Мащабиране от 1 до 6 реплики с цел на натоварването 0{,}70. Клиентът е
\code{stratus-loadgen}, изпълнен от София през публичния интернет, по обикновен HTTP на порт 80.

\textbf{Тези числа не са сравними с локалните.} Клиентът е в София, услугата е в Милано, а между
тях стои публичният интернет. Закъснението съдържа пътя през глобалната мрежа и не бива да се
поставя до таблица~\ref{tab:throughput}, където генераторът и контейнерите са на една машина.

Базовият пробег е с една връзка, 30 s и без загряване. Вторият е с 32 едновременни връзки, 10 s
загряване и 240 s измерване, и отчита 1750{,}43 Miter/s извършена работа \tabref{tab:azure}.

\begin{table}[H]
\caption{Azure Container Apps, два пробега от един и същи клиент. Суровите извадки от втория са в
\code{results/azure\_scaling.csv}}
\label{tab:azure}
\centering
\small
\begin{tabular}{@{}lrr@{}}
\toprule
\textbf{Показател} & \textbf{Базов пробег} & \textbf{Пробег с мащабиране} \\
\midrule
Успешни заявки             & 66      & 5698     \\
Неуспешни заявки           & 0       & 47       \\
Пропускателна способност [зая./s] & 2{,}20 & 23{,}74 \\
Средно закъснение [ms]     & 592{,}61   & 1282{,}74 \\
p50 [ms]                   & 120{,}24   & 1038{,}13 \\
p90 [ms]                   & 1129{,}83  & 2655{,}77 \\
p95 [ms]                   & 1133{,}61  & 3445{,}10 \\
p99 [ms]                   & 15157{,}82 & 15160{,}97 \\
Най-голямо [ms]            & 15157{,}82 & 24152{,}75 \\
\bottomrule
\end{tabular}
\end{table}

\textbf{Неуспешните заявки.} 47 от 5745 заявки във втория пробег са неуспешни, тоест 0{,}82\%. Разпределението
не е равномерно: 29 от тях са между 30-ата и 60-ата секунда, докато цялото натоварване се поема от
една реплика. Останалите 18 са разпръснати до края.

\textbf{Опашката не идва от едновременността.} p99 е около 15 s и в двата пробега, включително в
базовия, където има само една връзка и 66 заявки. Стойност, която се появява еднакво при 1 и при 32
едновременни заявки, не се обяснява с натрупана опашка. Данните не позволяват избор между мрежата
и студен старт, затова се съобщава само наблюдението.

\begin{table}[H]
\caption{Брой реплики по време на пробега с мащабиране, секунди от началото на натоварването.
Пълната редица е в \code{results/azure\_measurement.txt}}
\label{tab:azure-replicas}
\centering
\small
\begin{tabular}{@{}rr@{\hspace{1.2cm}}rr@{\hspace{1.2cm}}rr@{}}
\toprule
\textbf{$t$ [s]} & \textbf{Репл.} & \textbf{$t$ [s]} & \textbf{Репл.} &
\textbf{$t$ [s]} & \textbf{Репл.} \\
\midrule
3   & 1 & 108 & 2 & 230 & 3 \\
60  & 1 & 155 & 2 & 249 & 3 \\
99  & 1 & 221 & 2 & 268 & 3 \\
\bottomrule
\end{tabular}
\end{table}

\textbf{Време за реакция.} Това е най-полезното от пробега \tabref{tab:azure-replicas}. Първото
увеличаване от 1 на 2 реплики е настъпило \emph{между 99 и 108 секунди} след началото на
натоварването. Второто, от 2 на 3 реплики, е настъпило \emph{между 221 и 230 секунди}. Границата
\code{max\_replicas} от 6 реплики не е достигната в рамките на 240-секундния прозорец. Съобщава се
интервал, а не момент: броят реплики се чете на около 9 до 10 секунди, така че действителната смяна
е някъде между две последователни отчитания. Значение има порядъкът: реакцията е от порядъка на
минута и половина, докато локалният контролер реагира за секунди \tabref{tab:timeline}.

\subsection{Какво разкри разгръщането}

\textbf{Дефект в самото описание на инфраструктурата.} Във файла \code{infra/versions.tf} стоеше
коментар, че двата доставчика могат да се конфигурират безусловно, защото „нито един не се
удостоверява, преди действително да се планира ресурс, така че конфигурирането на неизползвания не
струва нищо“. Това е невярно. Доставчикът за AWS търси данни за достъп при
конфигурирането, а не при планирането на ресурс. Затова \code{terraform plan} за целта Azure изчисляваше пълния
план от девет ресурса за Azure и после се проваляше поради липсващи данни за достъп до AWS, при
това след изчакване на точката за метаданни на EC2 на адрес \code{169.254.169.254}. Грешката е
стояла незабелязана, защото \code{terraform} никога не е бил изпълняван. Поправена е с трите документирани аргумента \code{skip\_*} и защитни фиктивни ключове,
всички обусловени от \code{target\_provider}, така че действителната верига за удостоверяване
остава непокътната, когато целта е AWS.

\textbf{Наемателят забранява удостоверяване със споделен ключ за хранилищата.} Прилагането се
проваляше с \code{KeyBasedAuthenticationNotPermitted}, след като сметката за хранилище вече е
създадена. Модулът беше правилен: в него \code{shared\_access\_key\_enabled} е \code{false}. Вината е у доставчика, който по
подразбиране проверява равнището за данни със споделен ключ. Поправено с \code{storage\_use\_azuread = true} на доставчика \code{azurerm}.

\textbf{Azure for Students ограничава регионите.} Правило с име \code{sys.regionrestriction}
разрешава само \code{francecentral}, \code{switzerlandnorth}, \code{spaincentral},
\code{italynorth} и \code{denmarkeast}. Първият опит, в \code{polandcentral}, беше отказан.
Избран е \code{italynorth}, най-близкият разрешен регион до София.

\subsection{Разход и експлоатационни бележки}

\textbf{Разходът се вижда, но не се разпределя.} Наличният кредит спадна от 100 на 88 USD за около един
час, през който разгръщането беше вдигнато, докато полето „разходи за август“ в портала показваше
0{,}00. Тоест 12 USD са изразходвани и се виждат в остатъка от кредита,
но не са отнесени към услуга. При това темпо кредитът стига за часове, а не за 8760-те часа в
годината, за които изглежда предвиден. Затова всичко е унищожено веднага след пробега.

\textbf{Ингресът беше превключен ръчно.} Изпълнена е командата
\code{az containerapp ingress update --allow-insecure}, извън \code{terraform} и нарочно.
Генераторът на натоварване в този проект говори обикновен HTTP, а ингресът иначе отговаря на порт
80 с пренасочване 301. Промяната е извън описанието на инфраструктурата, за да запази то
поведението си само по HTTPS.

\subsection{Какво не е измерено}

\textbf{Няма сравнение между двама доставчици.} Това е най-тежкото ограничение на документа. При
AWS не е разгръщано, защото за тази сметка няма кредит. Замисълът, едно описание при двама доставчици и
измерена разлика между тях, остава проверен наполовина.

\textbf{Цената за единица работа остава неизмерена}, и то по причина, която не се решава с още
време. Програмният интерфейс Cost Management отказва тази абонаментна сметка направо, с код
HTTP 422: той поддържа само типовете Enterprise Agreement, Web direct и Microsoft Customer
Agreement, а Azure for Students не е нито един от трите. Остава само порталът, който показва
обобщена стойност на кредита и изостава спрямо потреблението. Тоест числото не е просто
незаписано, а е недостъпно през интерфейса при този тип абонамент.
